\documentclass{article}
\usepackage{amsmath,amsthm,amssymb}

\DeclareMathOperator{\area}{area}
\DeclareMathOperator{\dinv}{dinv}

\title{Dyck Skeleton Tableau Formulas}
\author{}
\date{}

\begin{document}
\maketitle

\section{What the formula says}

The purpose of this item is to compare two ways of forming the same
two-variable generating function.  The direct side sums
\(q^{\area}t^{\dinv}\) over normalized rational Dyck sequences.  The
skeleton/tableau side groups the same expected polynomial by rational
skeletons, at-most-two-column rational Dyck tableaux, and a two-variable
Schur factor.

The checker uses the option name \(t\) for the rational step and \(n\) for the
length.  In this explanation, the rational step is written \(\tau\) and the
length is written \(s\), so the congruence family is
\[
  r=s\tau+1.
\]
Thus the checker's case \(t=1\) is the case \(\tau=1\).  From this point on,
\(\tau\) denotes the rational step, while \(t\) remains the second variable in
\(q,t\).

\subsection{Rational Dyck side}

Fix \(s\geq 1\) and \(\tau\geq 0\).  A normalized rational Dyck sequence of
length \(s\) and step \(\tau\) is a sequence
\[
  D=(D_0,\ldots,D_{s-1})
\]
of nonnegative integers such that \(D_0=0\) and
\[
  D_{i+1}\leq D_i+\tau
  \qquad(0\leq i<s-1).
\]
For a finite integer sequence \(x=(x_0,\ldots,x_{\ell-1})\), define
\[
  \dinv_\tau(x)=
  \sum_{0\leq i<j<\ell} d_\tau(x_i,x_j),
\]
where
\[
  d_\tau(a,b)=
  \begin{cases}
    \max(0,a+\tau-b),& a\leq b,\\
    \max(0,b+1+\tau-a),& a>b.
  \end{cases}
\]
The area statistic is
\[
  \area(x)=\sum_i x_i.
\]
The direct rational \(q,t\)-Catalan side in this model is therefore
\[
  C_{s,\tau}(q,t)
  =
  \sum_D q^{\area(D)}t^{\dinv_\tau(D)},
\]
where \(D\) ranges over all normalized rational Dyck sequences of length
\(s\).

\subsection{Skeleton--tableau side}

The formula decomposes the same polynomial by first choosing a rational
skeleton and then filling the remaining cells by a two-column rational Dyck
tableau.

An entry \(D_j=e\) of a normalized rational Dyck sequence is called
extractable if \(e>0\), exactly one earlier entry lies in the predecessor
window
\[
  \{a\in\mathbb Z_{\geq 0}: \max(0,e-\tau)\leq a\leq e-1\},
\]
and deleting \(D_j\) preserves the adjacent rational Dyck inequality.  More
explicitly, if \(0<j<s-1\), then deletion requires
\[
  D_{j+1}\leq D_{j-1}+\tau.
\]
If \(j=s-1\), there is no new adjacent pair to check.  For \(\tau=0\), the
predecessor window is empty.  A rational \(m\)-skeleton is a normalized
rational Dyck sequence \(F\) with final entry equal to its maximum,
\[
  F_{|F|-1}=\max(F)=m,
\]
and with no nonfinal extractable entry.  A final extractable entry is allowed
in an \(m\)-skeleton.

A rational Dyck tableau \(P\) of step \(\tau\) is a left-aligned tableau whose
row lengths form a partition.  We index its rows in the source-paper
convention: \(P_0\) is the top row, \(P_1\) is the row immediately below it,
and so on.  Rows are read left to right and are dual rational Dyck sequences,
\[
  P_i[j+1]>P_i[j]+\tau,
\]
while columns, read bottom to top, satisfy the affine rational Dyck condition.
Equivalently, if row \(i\) is immediately above row \(i+1\), then
\[
  P_i[j]\leq P_{i+1}[j]+\tau
\]
whenever both entries exist.  Its row-reading word is
\[
  \operatorname{RR}(P)=P_{\text{bottom}}P_{\text{next}}\cdots P_{\text{top}},
\]
with each row read left to right, and \(\lambda(P)\) denotes its shape.

The checked conjectural rational two-column skeleton/tableau identity is
\[
  C_{s,\tau}(q,t)
  =
  \sum_{(F,P)}
  q^{\area(F:\operatorname{RR}(P))}
  t^{\dinv_\tau(F:\operatorname{RR}(P))-|P|}
  s_{\lambda(P)'}(q,t).
\]
The sum is over all pairs \((F,P)\) such that \(F\) is a rational
\(m\)-skeleton of step \(\tau\) for some \(m\geq 0\), \(P\) is an
at-most-two-column rational Dyck tableau of step \(\tau\) with entries in
\([0,m-1]\), and
\[
  |F|+|\operatorname{RR}(P)|=s.
\]
Here \(F:\operatorname{RR}(P)\) means concatenation, \(|P|\) is the number of
cells of \(P\), \(\lambda(P)'\) is the conjugate partition, and
\(s_{\lambda(P)'}(q,t)\) is the Schur function in the two variables \(q,t\).
This is the checked conjectural identity for general rational step
\(\tau\).

\section{Source context and status}

Status summary: the identity is proved for \(\tau=1\), degenerate for
\(\tau=0\), and computationally verified but not proved for the tested
values \(\tau>1\).

For \(\tau=1\), the formula is exactly the two-column tableau formula for the
ordinary \(q,t\)-Catalan polynomial proved in the source paper.  In that
proof, Dyck sequences are sent by the paper's Type 1 through Type 4
correspondences to Type 4 triples \((F,P,Q)\), where \(F\) is a Dyck
\(m\)-skeleton, \(P\) is an at-most-two-column Dyck tableau, and \(Q\) is a
binary reverse semistandard tableau of shape \(\lambda(P)\).

The \(-|P|\) shift in the displayed formula comes from summing over the
binary recording tableau \(Q\).  For fixed \((F,P)\), one has
\[
  q^{\#1(Q)}t^{-\#1(Q)}
  =
  t^{-|P|}q^{\#1(Q)}t^{\#0(Q)},
\]
since \(\#0(Q)+\#1(Q)=|P|\).  Transposing \(Q\) turns the binary reverse
semistandard condition on shape \(\lambda(P)\) into the ordinary
semistandard condition on the conjugate shape \(\lambda(P)'\), giving
\[
  \sum_Q q^{\#1(Q)}t^{\#0(Q)}
  =
  s_{\lambda(P)'}(t,q)
  =
  s_{\lambda(P)'}(q,t).
\]

The case \(\tau=0\) is degenerate.  The adjacent condition becomes
\(D_{i+1}\leq D_i\), so a normalized nonnegative sequence starts at \(0\) and
can only stay at \(0\).  Thus the direct side has only the all-zero sequence
in each length.  On the formula side, a contributing skeleton must have
\(m=0\), and a nonempty tableau would need entries in \([0,m-1]=[0,-1]\),
which is impossible.  Therefore only \(F=(0,\ldots,0)\) with \(P=\varnothing\)
contributes.

For \(\tau>1\), this item records finite evidence for the conjectural
rational analogue tested here.  The checker does not prove the formula in
these cases; it enumerates both sides and compares the resulting coefficient
dictionaries grouped by \((\area,\dinv)\).

\section{Computational verification}

The checker for the rational two-column skeleton/tableau formula is recorded
in
\[
\texttt{code/check\_rational\_two\_column\_formula.py}.
\]
It compares two coefficient dictionaries: one from the direct normalized
rational Dyck sequence generating function, and one from the formula-side
expansion described above.  The check passes exactly when the two
dictionaries agree for every \((\operatorname{area},\operatorname{dinv})\)
key in the requested finite range.

The official finite checks performed for this item are:
\begin{itemize}
\item \(\tau=2\), lengths \(1 \leq s \leq 14\), elapsed time
      \(4102.814\) seconds;
\item \(\tau=3\), lengths \(1 \leq s \leq 12\), elapsed time
      \(25254.334\) seconds;
\item \(\tau=4\), lengths \(1 \leq s \leq 10\), elapsed time
      \(494.131\) seconds.
\end{itemize}
All three official checks passed.  These finite checks do not prove the
identity for all lengths for any \(\tau>1\).

\end{document}
